\section*{الفصل \ref{ch:g3:sorting-living-things} --- فرز الكائنات الحية}

\begin{solution}{exo:g3:sorting-living-things:1}
\omterm{def:g3:sorting-living-things:criterion}{المحكّ} هو القاعدة المستعملة في \omterm{def:g3:sorting-living-things:criterion}{الفرز}. والجيد منه: صفة في الجسم
تستطيع ملاحظتها --- ``له \omterm{def:g1:animals-around-us:covering}{ريش}''، ``له ست \omterm{def:g1:my-body:parts}{أرجل}''. والمتجنَّب: أين
يعيش أو ما يصادف أن يفعله --- ``يعيش قرب الماء''، ``يطير''.
\end{solution}

\begin{solution}{exo:g3:sorting-living-things:2}
\omterm{prop:g3:sorting-living-things:groups}{الحشرات}: ست \omterm{def:g1:my-body:parts}{أرجل} وجسم من ثلاثة أقسام. \omterm{prop:g3:sorting-living-things:groups}{والأسماك}: زعانف وحراشف،
وحياة في الماء. \omterm{prop:g3:sorting-living-things:groups}{والطيور}: \omterm{def:g1:animals-around-us:covering}{ريش} وأجنحة ومنقار. \omterm{prop:g3:sorting-living-things:groups}{والثدييات}: \omterm{def:g1:animals-around-us:covering}{فرو} أو
شعر، وأمهات تُرضع صغارها \omterm{def:g2:animals-grow-up:born}{الحليب}.
\end{solution}

\begin{solution}{exo:g3:sorting-living-things:3}
النحلة: حشرة. والسمكة الذهبية: سمكة. والنعامة: طائر. والبقرة:
ثديي. والضفدع: برمائي.
\end{solution}

\begin{solution}{exo:g3:sorting-living-things:4}
له \omterm{def:g1:animals-around-us:covering}{فرو} وهو يُرضع صغاره \omterm{def:g2:animals-grow-up:born}{الحليب}، وليس له \omterm{def:g1:animals-around-us:covering}{ريش} ولا منقار. والطيران شيء
يفعله لا شيء يحدّ الطائر --- والصفات تقول: ثديي.
\end{solution}

\begin{solution}{exo:g3:sorting-living-things:5}
لا حراشف له ولا خطة جسم السمكة ذات الغلاصم والزعانف --- وهو يُرضع
عجله \omterm{def:g2:animals-grow-up:born}{الحليب}. والسباحة عادته؛ أما صفاته فتقول: ثديي.
\end{solution}

\begin{solution}{exo:g3:sorting-living-things:6}
لا. فمحكّ \omterm{prop:g3:sorting-living-things:groups}{الحشرات} ست \omterm{def:g1:my-body:parts}{أرجل} (وجسم من ثلاثة أقسام)؛ \omterm{def:g1:my-body:parts}{وأرجل} العنكبوت
الثماني تسقطه. والعنكبوت من مجموعة أخرى، مهما بدا لأول وهلة.
\end{solution}

\begin{solution}{exo:g3:sorting-living-things:7}
الأقحوان: أزهار \omterm{def:g2:from-seed-to-plant:seed}{وبذور}، \omterm{def:g1:plants-around-us:parts}{وساق} طرية خضراء. والطحلب: لا أزهار له
البتة ولا \omterm{def:g1:plants-around-us:parts}{ساق} حقيقية. والصنوبر: شجرة (جذع خشبي)، وأوراقه إبر، ولا
أزهار له بل أكواز فيها \omterm{def:g2:from-seed-to-plant:seed}{بذور}. وشجيرة الورد: أزهار \omterm{def:g2:from-seed-to-plant:seed}{وبذور}، \omterm{def:g1:plants-around-us:parts}{وسيقان}
خشبية.
\end{solution}

\begin{solution}{exo:g3:sorting-living-things:8}
الجواب مفتوح. وفي أكثر الشوارع تفوز \omterm{prop:g3:sorting-living-things:groups}{الحشرات} بفارق كبير --- نمل،
وذباب، ونحل، وخنافس، وفراشات --- \omterm{prop:g3:sorting-living-things:groups}{والطيور} في المرتبة الثانية.
\end{solution}

\begin{solution}{exo:g3:sorting-living-things:9}
\omterm{def:g3:sorting-living-things:criterion}{فرز} الصديقة. \omterm{def:g3:sorting-living-things:criterion}{فالفرز} إلى ``أليفة'' و``برية'' يجري بحسب طريقة عيش
\omterm{def:g1:animals-around-us:animal}{الحيوانات} مع الناس، وهي قابلة للتغير --- فالقطة الشاردة هي \omterm{def:g1:animals-around-us:animal}{الحيوان}
نفسه داخل البيت وخارجه. أما ``\omterm{prop:g3:sorting-living-things:groups}{الثدييات} في مقابل \omterm{prop:g3:sorting-living-things:groups}{الطيور}'' فيفرز
بصفات جسم تبقى مدى الحياة: وهو محكّ عالم الأحياء.
\end{solution}

\begin{solution}{exo:g3:sorting-living-things:10}
الخطوة الأولى: الكساء --- \omterm{def:g1:animals-around-us:covering}{ريش} (متراصّ)، ومنقار، \omterm{def:g1:my-body:parts}{ورجلان}، وجناحان
(على هيئة زعنفتين). والخطوة الثانية: الصغار تفقس من \omterm{def:g2:animals-grow-up:born}{بيض}. والخطوة
الثالثة: الصفات توافق \omterm{prop:g3:sorting-living-things:groups}{الطيور}. والخطوة الرابعة: الانطباعان ``يسبح
كالسمكة، ولا يطير'' يخالفان --- فيخسران. البطريق طائر.
\end{solution}
